\section{Curry--Howard: proofs are terms}

What is Curry Howard? A bit of history. What problem is it solving? How does it work? Why does it work?

What does it mean for a term to be the same as a proof and a proposition to be the same as a type? Why can we "understand proofs as terms that inhabit the type of the thing we are proving"? Why does that work?

\subsection{Proving things}

Ok so proving something is constructing a term. How do you even "construct a term"? Well remember the construction rules that we had at chapter 1.

Some bridge here im not really seeing. Like at this point i should have some tool box that allows me to build very simple proofs.

Begin some examples about proving things in term mode.

For example

prove P -> P 


So what we want to do is construct a term of type P -> P. It is simple, we can construct t with what we know.

Notice that if t is to have type P -> P, then we understand t as "a function from type P to type P", i.e., a function that takes elements from the type P and produces elements from the type P. An element from the type P is actually a proof of P as we saw. So from an arbitrary term of type P, say p : P, we need to construct a term of type P. Well, p is already of type P so simply return p.

I.e. we want to construct the function

$$
\begin{array}{ccrcl}
  t & : & P  &\longrightarrow  & P \\
  & & p & \longmapsto & p
\end{array}
$$

What is this function in lambda terms? It is simply $\lambda p : P . p$.

More examples.

\subsubsection{First order logic}

Deduction rules for first order logic. This belongs here somehow but I am unsure how. 

Like curry howard also defines a way to turn NJ trees into lambda terms no?

So this way we could examples like for example A and B <-> B and A

\subsubsection{Exercises}

\subsection{Proving more difficult things}

Tactic mode, the elaborator. Understanding that a proof is a set of steps i.e. a proof tree.

\subsubsection{Exercises}

